\documentclass[12pt, a4paper]{article}
\usepackage[utf8]{inputenc}
\usepackage[spanish]{babel}
\usepackage{hyperref}
\usepackage{geometry}
\usepackage{titlesec}
\usepackage{enumitem}
\usepackage{tabularx}
\usepackage{booktabs}
\usepackage{parskip}

\geometry{a4paper, margin=1in}

\title{\textbf{Especificación de Requisitos de Software (SRS)}\\ \Large Comercial de Mariscos Don Walter (PescaderiaLaravel)}
\author{Documento basado en el estándar IEEE Std 830-1998}
\date{\today}

\begin{document}

\maketitle
\tableofcontents
\newpage

\section{Introducción}

\subsection{Propósito}
El propósito de este documento es definir la Especificación de Requisitos de Software (SRS) para el sistema de e-commerce y gestión de la ``Comercial de Mariscos Don Walter'', desarrollado bajo el framework Laravel 9. Este documento detalla tanto las funcionalidades del frontend (tienda orientada al cliente) como del backend (panel CRM de administración).

\subsection{Ámbito del Sistema}
El software es una plataforma integral de comercio electrónico para pescado y mariscos de alta calidad. El sistema permite a los clientes explorar productos mediante una interfaz inmersiva, gestionar un carrito de compras y realizar pedidos. Simultáneamente, ofrece a la administración un Dashboard de gestión de inventarios, despachos, métricas financieras en tiempo real y configuraciones estrictas de seguridad (Hardening y rotación de contraseñas).

\subsection{Definiciones, Acrónimos y Abreviaturas}
\begin{itemize}
    \item \textbf{SRS}: Software Requirements Specification.
    \item \textbf{PDP}: Product Details Page (Página de Detalles de Producto).
    \item \textbf{CRM}: Customer Relationship Management.
    \item \textbf{ORM}: Object-Relational Mapping (Eloquent en Laravel).
    \item \textbf{RUC}: Registro Único de Contribuyentes.
\end{itemize}

\subsection{Referencias}
\begin{itemize}
    \item IEEE Std 830-1998, IEEE Recommended Practice for Software Requirements Specifications.
    \item Documentación oficial de Laravel 9.0 (\url{https://laravel.com/docs/9.x}).
\end{itemize}

\subsection{Visión General del Documento}
El presente documento está estructurado de acuerdo con el estándar IEEE 830-1998. La sección 2 ofrece una perspectiva general del sistema. La sección 3 describe exhaustivamente los requisitos funcionales, no funcionales y las interfaces externas.

\section{Descripción General}

\subsection{Perspectiva del Producto}
El sistema es una aplicación web autónoma que reemplaza los métodos tradicionales de venta física e inventario manual. Está diseñado sobre una arquitectura Modelo-Vista-Controlador (MVC) conectada a una base de datos relacional (MySQL).

\subsection{Funciones del Producto}
Las funciones principales del sistema incluyen:
\begin{itemize}
    \item \textbf{Gestión de Catálogo:} Operaciones CRUD sin refresco de página (AJAX) para Categorías y Productos.
    \item \textbf{Escaparate Digital (E-commerce):} Galería de productos, Página de Detalles Inmersiva (PDP) al estilo Amazon y carrito de compras dinámico.
    \item \textbf{Gestión de Órdenes:} Creación nativa de pedidos por parte de clientes y módulo de despachos y facturación (SRI) en el panel administrativo.
    \item \textbf{Panel Analítico:} Generación de gráficos Chart.js con balances financieros y existencias.
    \item \textbf{Módulo de Seguridad:} Configuración del sistema y notificaciones proactivas forzando la rotación de credenciales maestras cada 90 días.
\end{itemize}

\subsection{Características de los Usuarios}
Existen 2 perfiles principales interactuando con la aplicación:
\begin{enumerate}
    \item \textbf{Clientes:} Consumidores finales de servicios web. Requieren una interfaz sumamente intuitiva y atractiva visualmente sin requerir conocimientos técnicos.
    \item \textbf{Administradores:} Gestores logísticos y contables con requerimientos de alta disponibilidad sobre tableros de inventario y manipulación de datos sensibles.
\end{enumerate}

\subsection{Entorno Operativo}
\begin{itemize}
    \item Servidor web: Mínimo Apache 2.4 o Nginx.
    \item Lenguaje base: PHP 8.0+.
    \item Base de Datos: MySQL 5.7+ o MariaDB.
    \item Frameworks UI: Bootstrap 5.3, DataTables.
\end{itemize}

\subsection{Restricciones de Diseño e Implementación}
\begin{itemize}
    \item La aplicación debe hacer uso exhaustivo de PDO Parameter Binding en Eloquent para mitigar inyecciones SQL.
    \item Las vistas deben adherirse estricamente a la metodología \textit{Mobile-first} mediante el sistema grid de Bootstrap.
    \item El borrado de categorías no debe eliminar en cascada los productos en la base de datos (Restricción Set Null).
\end{itemize}

\section{Requisitos Específicos}

\subsection{Requisitos de Interfaz Externa}
\subsubsection{Interfaces de Usuario}
El sistema presentará una separación de interfaces estricta entre Administradores y Clientes. La interfaz Frontal debe incluir un \textit{Buy Box} y compatibilidad con pantallas táctiles, mientras que el \textit{Backend Dashboard} alojará paneles colapsables laterales y componentes DataTables.

\subsubsection{Interfaces de Hardware y Software}
La plataforma deberá ejecutarse eficazmente sobre cualquier navegador web moderno (Chrome, Firefox, Safari, Edge). Carece de interfaces de integración al hardware más allá del soporte de impresión para los reportes o comprobantes en formato A4.

\subsection{Requisitos Funcionales}

\subsubsection{RF-01: Autenticación de Usuarios}
El sistema debe proveer vistas nativas e independientes para el Registro de usuarios y el Inicio de sesión. La verificación del usuario será contra un correo y contraseña cifrada usando Bcrypt. Las validaciones regresarán un HTTP 422 JSON explícito en casos de fallos en requisitos formales.

\subsubsection{RF-02: Exploración de Productos E-commerce}
Los clientes podrán explorar los productos y ser redirigidos a una PDP (Product Details Page) de 3 columnas divididas visualmente en: Galería (Col 1), Ficha técnica (Col 2) y Caja de Compra (Col 3) calculada por unidad de medida (\textit{lbs, un}).

\subsubsection{RF-03: Procesamiento del Carrito de Compras}
El cliente debe poder incrementar las unidades asíncronamente o dentro del entorno sincrónico validando en todo momento que su solicitud no supere la existencia (\textit{stockVenta}) en la base de datos local.

\subsubsection{RF-04: Panel Estadístico de Despachos}
El aplicativo facilitará a los administradores gráficos estadísticos renderizados con Chart.js basándose en las órdenes consolidadas semanal y mensualmente.

\subsubsection{RF-05: Modulación de Seguridad Activa}
La plataforma auditará recurrentemente a los administradores. Si las credenciales concuerdan con los \textit{default factory hashes} (admin123) o si superan los 90 días en latencia sin un Hash update, se inyectarán Modales o Alertas interrumpiendo el flujo logístico hasta que la contraseña se restablezca.

\subsection{Requisitos de Rendimiento}
\begin{itemize}
    \item \textbf{Tiempos de respuesta:} El sistema deberá responder a interacciones E-commerce (Add to Cart) en un lapso no superior a los 500ms al servidor bajo demanda transaccional promedio.
    \item \textbf{Simultaneidad:} La arquitectura MVC y base de datos relacional soportarán múltiples operadores de despachos modificando existencias a la par, apalancándose de transacciones SQL.
\end{itemize}

\subsection{Atributos del Sistema}
\subsubsection{Seguridad}
- Uso exclusivo de directivas del ORM para evitar fallos de inyección SQL. \\
- Protección CSRF adherida globalmente en todos los formularios HTML de método POST/PUT/DELETE. \\
- Segmentación por middlewares en rutas administrativas y de clientes.

\subsubsection{Fiabilidad y Mantenibilidad}
La estructura de directorios sigue estrictamente el ecosistema formal de Laravel. El desacoplamiento entre las \textit{Controllers} y los modelos asegura un sencillo marco para escalar horizontalmente la plataforma.

\end{document}
